% GENERATED FILE. Edit canonical lessons, then run npm run generate:books.
% canonical-source-hash: fnv1a64:52e33cd1cb19b6a3
% canonical-lessons: JA-C408-senkou, JA-C408-rousoku, JA-C408-butsudan, JA-C408-shuukyou, JA-C408-matsuru

\chapter{Incense and Candles}
\label{ch:ja-incense-and-candles}
\hlchaptermodality{408}

\begin{chapteropening}
\textbf{By the end of this chapter:} \emph{I can talk about prayer, altars and religion.}

Complete the last lesson of chapter 408: I can talk about prayer, altars and religion.
\end{chapteropening}

\section[\texorpdfstring{\ja{せんこう} (senkō)}{せんこう (senkō)}]{\ja{せんこう} (senkō) — an incense stick}

\label{lesson:JA-C408-senkou}



\begin{quote}
Before the new one: say the Japanese for the main hall of a temple, then the Japanese for to pray with joined hands.
\end{quote}

\subsection*{You'll want to know: \ja{せんこう}}
\textbf{\ja{せんこう}} — \emph{senkō} — \textquotedblleft{}an incense stick\textquotedblright{}.

\begin{quote}
\textbf{\ja{ほんどうの} \ja{まえで} \ja{せんこうを} \ja{たてました。}} — \emph{hondō no mae de senkō o tatemashita} — \textquotedblleft{}We stood incense sticks in front of the main hall.\textquotedblright{}
\end{quote}

\begin{cousinweb}
\textbf{\ja{せんこう}} (\emph{senkō}) is Sino-Japanese (kango): two Chinese characters read with their on'yomi, readings borrowed from Chinese. They are \emph{sen}, \textquotedblleft{}line\textquotedblright{} (Middle Chinese \emph{sjenH}) and \emph{kō}, \textquotedblleft{}incense\textquotedblright{} (\emph{xjang}) (KANJIDIC2; Middle Chinese in Baxter's transcription).
\end{cousinweb}

\begin{grammarlens}[title={What you've built}]
One more word for getting around.
\end{grammarlens}

\subsection*{Your turn}
\begin{itemize}
\raggedright
  \item \emph{Say it:} \emph{senkō}
  \item \emph{Say it:} \emph{senkō}, once more
  \item \emph{Say it:} \emph{hondō no mae de senkō o tatemashita} — We stood incense sticks in front of the main hall.
  \item \emph{Recall:} say the Japanese for the main hall of a temple, then the Japanese for to pray with joined hands, then say \emph{senkō} again
\end{itemize}

\begin{tcolorbox}[breakable,colback=teal!4,colframe=teal!35!black,title={Before you move on}]
What does \ja{せんこう} mean? (\textquotedblleft{}An incense stick\textquotedblright{}.) Say it once more.
\end{tcolorbox}
\section[\texorpdfstring{\ja{ろうそく} (rōsoku)}{ろうそく (rōsoku)}]{\ja{ろうそく} (rōsoku) — a candle}

\label{lesson:JA-C408-rousoku}



\begin{quote}
Before the new one: say the Japanese for to pray with joined hands, then the Japanese for an incense stick.
\end{quote}

\subsection*{You'll want to know: \ja{ろうそく}}
\textbf{\ja{ろうそく}} — \emph{rōsoku} — \textquotedblleft{}a candle\textquotedblright{}.

\begin{quote}
\textbf{\ja{ろうそくの} \ja{ひが} \ja{きえました。}} — \emph{rōsoku no hi ga kiemashita} — \textquotedblleft{}The candle went out.\textquotedblright{}
\end{quote}

\begin{cousinweb}
\textbf{\ja{ろうそく}} (\emph{rōsoku}) is Sino-Japanese (kango): two Chinese characters read with their on'yomi, readings borrowed from Chinese. They are \emph{rō}, \textquotedblleft{}wax\textquotedblright{} (Middle Chinese \emph{lap}) and \emph{soku}, \textquotedblleft{}light\textquotedblright{} (\emph{tsyowk}) (KANJIDIC2; Middle Chinese in Baxter's transcription).
\end{cousinweb}

\begin{grammarlens}[title={What you've built}]
One more word for getting around.
\end{grammarlens}

\subsection*{Your turn}
\begin{itemize}
\raggedright
  \item \emph{Say it:} \emph{rōsoku}
  \item \emph{Say it:} \emph{rōsoku}, once more
  \item \emph{Say it:} \emph{rōsoku no hi ga kiemashita} — The candle went out.
  \item \emph{Recall:} say the Japanese for to pray with joined hands, then the Japanese for an incense stick, then say \emph{rōsoku} again
\end{itemize}

\begin{tcolorbox}[breakable,colback=teal!4,colframe=teal!35!black,title={Before you move on}]
What does \ja{ろうそく} mean? (\textquotedblleft{}A candle\textquotedblright{}.) Say it once more.
\end{tcolorbox}
\section[\texorpdfstring{\ja{ぶつだん} (butsudan)}{ぶつだん (butsudan)}]{\ja{ぶつだん} (butsudan) — a household altar (for the Buddha)}

\label{lesson:JA-C408-butsudan}



\begin{quote}
Before the new one: say the Japanese for an incense stick, then the Japanese for a candle.
\end{quote}

\subsection*{You'll want to know: \ja{ぶつだん}}
\textbf{\ja{ぶつだん}} — \emph{butsudan} — \textquotedblleft{}a household altar (for the Buddha)\textquotedblright{}.

\begin{quote}
\textbf{\ja{そぼの} \ja{うちには} \ja{ぶつだんが} \ja{あります。}} — \emph{sobo no uchi ni wa butsudan ga arimasu} — \textquotedblleft{}There is a family altar in my grandmother's house.\textquotedblright{}
\end{quote}

\begin{cousinweb}
\textbf{\ja{ぶつだん}} (\emph{butsudan}) is Sino-Japanese (kango): two Chinese characters read with their on'yomi, readings borrowed from Chinese. They are \emph{butsu}, \textquotedblleft{}Buddha\textquotedblright{} (Middle Chinese \emph{bjut}) and \emph{dan}, \textquotedblleft{}podium\textquotedblright{} (\emph{dan}) (KANJIDIC2; Middle Chinese in Baxter's transcription).
\end{cousinweb}

\begin{grammarlens}[title={What you've built}]
One more word for getting around.
\end{grammarlens}

\subsection*{Your turn}
\begin{itemize}
\raggedright
  \item \emph{Say it:} \emph{butsudan}
  \item \emph{Say it:} \emph{butsudan}, once more
  \item \emph{Say it:} \emph{sobo no uchi ni wa butsudan ga arimasu} — There is a family altar in my grandmother's house.
  \item \emph{Recall:} say the Japanese for an incense stick, then the Japanese for a candle, then say \emph{butsudan} again
\end{itemize}

\begin{tcolorbox}[breakable,colback=teal!4,colframe=teal!35!black,title={Before you move on}]
What does \ja{ぶつだん} mean? (\textquotedblleft{}A household altar (for the Buddha)\textquotedblright{}.) Say it once more.
\end{tcolorbox}
\section[\texorpdfstring{\ja{しゅうきょう} (shūkyō)}{しゅうきょう (shūkyō)}]{\ja{しゅうきょう} (shūkyō) — religion}

\label{lesson:JA-C408-shuukyou}



\begin{quote}
Before the new one: say the Japanese for a candle, then the Japanese for a household altar.
\end{quote}

\subsection*{You'll want to know: \ja{しゅうきょう}}
\textbf{\ja{しゅうきょう}} — \emph{shūkyō} — \textquotedblleft{}religion\textquotedblright{}.

\begin{quote}
\textbf{\ja{せんせいと} \ja{しゅうきょうの} \ja{ことを} \ja{はなしました。}} — \emph{sensei to shūkyō no koto o hanashimashita} — \textquotedblleft{}I talked with my teacher about religion.\textquotedblright{}
\end{quote}

\begin{cousinweb}
\textbf{\ja{しゅうきょう}} (\emph{shūkyō}) is Sino-Japanese (kango): two Chinese characters read with their on'yomi, readings borrowed from Chinese. They are \emph{shū}, \textquotedblleft{}religion\textquotedblright{} (Middle Chinese \emph{tsowng}) and \emph{kyō}, \textquotedblleft{}teach\textquotedblright{} (\emph{kaewH}) (KANJIDIC2; Middle Chinese in Baxter's transcription).
\end{cousinweb}

\begin{grammarlens}[title={What you've built}]
One more word for getting around.
\end{grammarlens}

\subsection*{Your turn}
\begin{itemize}
\raggedright
  \item \emph{Say it:} \emph{shūkyō}
  \item \emph{Say it:} \emph{shūkyō}, once more
  \item \emph{Say it:} \emph{sensei to shūkyō no koto o hanashimashita} — I talked with my teacher about religion.
  \item \emph{Recall:} say the Japanese for a candle, then the Japanese for a household altar, then say \emph{shūkyō} again
\end{itemize}

\begin{tcolorbox}[breakable,colback=teal!4,colframe=teal!35!black,title={Before you move on}]
What does \ja{しゅうきょう} mean? (\textquotedblleft{}Religion\textquotedblright{}.) Say it once more.
\end{tcolorbox}
\section[\texorpdfstring{\ja{まつる} (matsuru)}{まつる (matsuru)}]{\ja{まつる} (matsuru) — to enshrine; to worship}

\label{lesson:JA-C408-matsuru}



\begin{quote}
Before the new one: say the Japanese for a household altar, then the Japanese for religion.
\end{quote}

\subsection*{You'll want to know: \ja{まつる}}
\textbf{\ja{まつる}} — \emph{matsuru} — \textquotedblleft{}to enshrine; to worship\textquotedblright{}.

\begin{quote}
\textbf{\ja{この} \ja{じんじゃは} \ja{やまを} \ja{まつっています。}} — \emph{kono jinja wa yama o matsutte imasu} — \textquotedblleft{}This shrine worships the mountain.\textquotedblright{}
\end{quote}

Polite \textbf{\ja{まつります}}, \emph{matsurimasu}; a godan verb. Past (polite): \textbf{\ja{まちの} \ja{みんなが} \ja{やまを} \ja{まつりました。}}, \emph{machi no minna ga yama o matsurimashita}.

\begin{cousinweb}
\textbf{\ja{まつる}} (\emph{matsuru}) is a native Japanese word (wago), inherited rather than borrowed from Chinese; its kanji was attached later as a kun'yomi (JMdict). It is the verb behind \emph{matsuri}, \textquotedblleft{}a festival\textquotedblright{} (JMdict).
\end{cousinweb}

\begin{grammarlens}[title={What you've built}]
One more word for getting around.
\end{grammarlens}

\subsection*{Your turn}
\begin{itemize}
\raggedright
  \item \emph{Say it:} \emph{matsuru}
  \item \emph{Say it:} \emph{matsuru}, once more
  \item \emph{Say it:} \emph{kono jinja wa yama o matsutte imasu} — This shrine worships the mountain.
  \item \emph{Recall:} say the Japanese for a household altar, then the Japanese for religion, then say \emph{matsuru} again
\end{itemize}

\begin{tcolorbox}[breakable,colback=teal!4,colframe=teal!35!black,title={Before you move on}]
What does \ja{まつる} mean? (\textquotedblleft{}To enshrine; to worship\textquotedblright{}.) Say it once more.
\end{tcolorbox}
